\documentclass[10pt,a4paper]{amsart}
\usepackage[margin=19mm]{geometry}
\usepackage{amsmath,amssymb,mathtools}
\usepackage[scr=rsfs,cal=euler]{mathalfa}
\usepackage{xfrac}
\usepackage[unicode,hidelinks,bookmarks=false]{hyperref}
\newcommand{\ie}{i.e.\ }
\newcommand{\cf}{cf.\ }
\newcommand{\resp}{respectively\ }
\newcommand{\Subsection}{\subsection}
\providecommand{\nobreakdash}{\nobreak\hbox{-}}
\setlength{\emergencystretch}{2em}
\multlinegap0pt
\usepackage{bm}
\usepackage{bbm}
\usepackage{dsfont}
\usepackage{xspace}
\usepackage{cancel}
\usepackage{mathtools}
\usepackage{amsthm}
\makeatletter
\@ifundefined{Theorem}{\newtheorem{Theorem}{Theorem}}{}
\makeatother
\newcommand {\diag}  {\mathop{\rm diag}\nolimits}
\newcommand {\mat} [1] {\left[\begin{array}{#1}}
\newcommand {\rix}     {\end{array}\right]}
\newcommand {\setC}  {{\mathbb C}}
\newcommand {\wt} {\widetilde}
\title[Combined perturbation bounds for eigenstructure of Hermitian]{Combined perturbation bounds for eigenstructure of Hermitian matrices and singular structure of general matrices}
\author{Xiao Shan Chen, Hongguo Xu}
\date{Source: arXiv:2509.10688v1 (CC BY 4.0)}
\begin{document}
\maketitle

\noindent\emph{Source and licence.} Combined perturbation bounds for eigenstructure of Hermitian matrices and singular structure of general matrices. Version arXiv:2509.10688v1 (CC BY 4.0), distributed under the Creative Commons Attribution 4.0 International licence, as recorded in the arXiv licence metadata for this exact version and shown on the abstract page https://arxiv.org/abs/2509.10688v1. Full rights notice: licenses/arxiv-2509.10688-CC-BY-4.0.txt.\par\smallskip

\noindent\textit{Editorial scope.} Counted unit: the Theorem whose source label is \texttt{them1} in the article (source0.tex lines 435--464, source label \texttt{them1}), delivered together with the article's own complete original proof of it (source0.tex lines 466--560, one \texttt{proof} environment of 95 lines). No mathematics is written, repaired or re-derived: statement and proof are the article's own text line for line. Required context retained uncounted: c-121 (lines 300--302); spectral (lines 381--385); bku (lines 387--389); bkl (lines 392--395); c-2-3 (lines 406--408). Every definition, display and label of those lines is retained unchanged. Omitted from this package and not redistributed: 1--299, 303--380, 386--386, 390--391, 396--405, 409--434, 465--465, 561--1326. Nothing inside a retained excerpt is omitted or altered: every retained body is delivered exactly as the article's lines are written, so no omission receipt of any kind accompanies this package. Citations: the retained excerpts carry no citation command at all, so no bibliography entry of the article is transcribed and no reference is redistributed. Reference adaptation: none was needed and none was made; every reference inside the delivered unit resolves inside it, so no unresolved reference or citation is produced. Printed slips kept literal: no printed word, symbol, hypothesis or number of the counted statement or of its proof was altered. Layout and class: the article's own class is \texttt{article} with packages including amssymb, amsmath, color, a4wide, xcolor, exscale, cmmib57, stmaryrd; the wrapper is the lane's pinned \texttt{amsart} editorial wrapper at 10pt on A4, which restates the theorem-like environments the retained text needs and adds the article's own macro definitions that the retained text needs (\texttt{\textbackslash diag}, \texttt{\textbackslash mat}, \texttt{\textbackslash rix}, \texttt{\textbackslash setC}, \texttt{\textbackslash wt}), with extra wrapper packages (bm, bbm, dsfont, xspace, cancel, mathtools). The delivered numbering is the unit's own. Assets: no retained span contains a figure environment, a graphics-inclusion command, a PDF-page inclusion, a table or a TikZ picture; no image file is copied into this package, the package declares no asset dependency and no image file is redistributed with it. Licence: version arXiv:2509.10688v1 of this article is distributed under the Creative Commons Attribution 4.0 International licence, as recorded in the arXiv licence metadata for this exact version and shown on the abstract page https://arxiv.org/abs/2509.10688v1. A full rights notice is delivered at licenses/arxiv-2509.10688-CC-BY-4.0.txt. Characters: every retained excerpt is pure ASCII, so the delivered engine, XeLaTeX with its default Latin Modern text font, reports no missing character.
\begin{eqnarray}\label{c-121}
\|{\rm Eig}^{\downarrow}(\wt A)-{\rm Eig}^{\downarrow}(A)\|_F\leq \|\Delta A\|_F.
\end{eqnarray}

\begin{eqnarray}\label{spectral}
A(t)=U(t)\Lambda(t)U(t)^H,
%=\left[U_1(t),U_2(t)\right]\mat{cc}\Lambda_1(t)&0\\0&\Lambda_2(t)\rix
%\mat{c}U_1(t)^H\\U_2(t)^H\rix,
\end{eqnarray}

\begin{eqnarray}\label{bku}
U(t)=\left[U_1(t),\ldots,U_k(t)\right],\qquad U_j(t)\in\setC^{n\times r_j}, \quad j=1,\ldots,k,
\end{eqnarray}

\begin{eqnarray}\label{bkl}
\Lambda(t)=\diag(\Lambda_1(t),\ldots,\Lambda_k(t)),\qquad
\Lambda_j(t)\in{\setC}^{r_j\times r_j},\quad j=1,\ldots,k
\end{eqnarray}

\begin{eqnarray}\label{c-2-3}
U_j(t)^H\dot U_j(t)=0,\quad j=1,\ldots,k.
\end{eqnarray}

\begin{Theorem}\label{them1}
Let $A(t)=A+t\Delta A\in\mathbb{C}^{n\times n}$ be a Hermitian matrix with $t\in{\mathbb R}$.
Suppose $A(t)$ has the analytic decomposition (\ref{spectral}) with $U(t)$ satisfying
(\ref{bku}) and (\ref{c-2-3}) and $\Lambda(t)$ in the block diagonal form (\ref{bkl}).
Define
\begin{eqnarray}\label{c-2-1}
\delta_{ji}(t)=\min_{\lambda_1(t)\in\lambda(\Lambda_j(t)),~
\lambda_2(t)\in\lambda(\Lambda_i(t))}\left\{\left|\lambda_1(t)-\lambda_2(t)\right|\right\},\quad
\delta_j(t)=\min_{i\ne j}\left\{\delta_{ji}(t)\right\},
\end{eqnarray}
and
\begin{eqnarray}\label{c-2-2}
\delta_{j,\min}=\min_{0\le t\le 1}\delta_j(t),
\end{eqnarray}
for $j=1,\ldots,k$. Denote $\wt A=: A(1)=A+\Delta A$ and
\begin{eqnarray*}
&&U(0)=\left[U_1(0),\ldots, U_k(0)]=:[U_1, \ldots, U_k\right]=U,\\
%&&\Lambda(0)=\diag(\Lambda_1(0),\ldots,\Lambda_k(0))=:\diag(\Lambda_1,\ldots,
%\Lambda_k)=\Lambda,\\
&&U(1)=\left[U_1(1),\ldots, U_k(1)\right]=:[\wt U_1,\ldots, \wt U_k ]=\wt U. %\\
%&&\Lambda(1)=\diag(\Lambda_1(1),\ldots,\Lambda_k(1))=: \diag(\wt\Lambda_1,\ldots,
%\wt\lambda_k)=\wt\Lambda.
\end{eqnarray*}
Then
\begin{eqnarray}\label{c-2-4}
\|{\rm Eig}^{\downarrow}(\wt A)-{\rm Eig}^{\downarrow}(A)\|_F^2+\sum_{j=1}^k
\delta_{j,\min}^2\|\wt U_j-U_j\|_F^2\le
\|{\Delta A}\|_F^2.
\end{eqnarray}
\end{Theorem}\vskip 2mm

\begin{proof}
By taking derivatives on both sides of (\ref{spectral}) and $U^H(t)U(t)=I,$ we have
\begin{eqnarray}\label{dspectral}
\Delta A=\dot{U}(t)\Lambda(t)U(t)^H+U(t)\dot{\Lambda}(t)U(t)^H+U(t)\Lambda(t)\dot{U}(t)^H
\end{eqnarray}
and
\begin{eqnarray}\label{dunitary}
\dot{U}(t)^H U(t)=-U(t)^H \dot{U}(t).
\end{eqnarray}
Multiplying both sides of (\ref{dspectral}) with $U(t)^H$ on the left and $U(t)$ on the right,
and using  (\ref{dunitary}), we obtain
\begin{eqnarray}\label{dspectral-2}
U(t)^H \Delta A U(t) =\dot{\Lambda}(t)+U(t)^H \dot{U}(t)\Lambda(t)-\Lambda(t) U(t)^H \dot{U}(t).
\end{eqnarray}
Let
\[
U_i(t)^H\Delta AU_j(t)=:[\Delta A_{ij}(t)],\quad i,j=1,\ldots,k.
\]
By comparing the blocks of (\ref{dspectral-2}) and using (\ref{c-2-3}), one obtain its diagonal terms
\begin{eqnarray}\label{c-2-5}
\dot\Lambda_j(t)=\Delta A_{jj}(t),\quad j=1,\ldots, k,
\end{eqnarray}
and off-diagonal terms
\begin{eqnarray}\label{c-2-6}
U_i(t)^H\dot U_j(t)\Lambda_j(t)-\Lambda_i(t)U_i(t)^H\dot U_j(t)=\Delta A_{ij}(t),
\qquad i\ne j.
\end{eqnarray}
Suppose that
\[
\Lambda_j(t)=G_j(t)\mat{ccc}\lambda_{j,1}(t)&&\\&\ddots&\\&&\lambda_{j,r_j}(t)\rix
G_j(t)^H,\quad j=1,\ldots,k,
\]
are  spectral decompositions (not necessarily analytic),
where $G_1(t),\ldots,G_k(t)$
are unitary matrices. Multiplying both sides of (\ref{c-2-6}) with $G_i(t)^H$ on the left and  $G_j(t)$ on the right yields
\begin{eqnarray}\label{c-2-7}
X_{ij}(t)\mat{ccc}\lambda_{j,1}(t)&&\\&\ddots&\\&&\lambda_{j,r_j}(t)\rix-
\mat{ccc}\lambda_{i,1}(t)&&\\&\ddots&\\&&\lambda_{i,r_i}(t)\rix X_{ij}(t)
=Z_{ij}(t),
\end{eqnarray}
where $X_{ij}(t)=G_i(t)^HU_i(t)^H\dot U_j(t)G_j(t)$
and $Z_{ij}(t)=G_i(t)^H\Delta A_{ij}(t)G_j(t)$. Let
$$X_{ij}(t)=[x^{(ij)}_{pq}(t)],\quad   Z_{ij}(t)=[z_{pq}^{(ij)}(t)].$$
 From (\ref{c-2-7}) one has
\begin{eqnarray*}
z^{(ij)}_{pq}(t)=(\lambda_{j,q}(t)-\lambda_{i,p}(t))x_{pq}^{(ij)}(t), \quad p=1,\ldots,r_i,\quad
 q=1,\ldots,r_j.
\end{eqnarray*}
Then for $\delta_{ji}(t)$ defined in (\ref{c-2-1}), one has
\begin{eqnarray}\label{c-2-8}
\|{\Delta A_{ij}(t)}\|_F^2&=&
\|{Z_{ij}(t)}\|_F^2=\sum_{p=1}^{r_i}\sum_{q=1}^{r_j}|z_{pq}^{(ij)}(t)|^2
=\sum_{p=1}^{r_i}\sum_{q=1}^{r_j}|\lambda_{j,q}(t)-\lambda_{i,p}(t)|^2|x_{pq}^{(ij)}(t)|^2
 \nonumber\\
&\ge& \delta_{ji}(t)^2\sum_{p=1}^{r_i}\sum_{q=1}^{r_j}|x_{pq}^{(ij)}(t)|^2
=\delta_{ji}(t)^2\|{U_i(t)^H\dot U_j(t)}\|_F^2.
\end{eqnarray}
By (\ref{c-2-3}),  we have
\begin{eqnarray}\label{c-2-9}
\|{\dot U_j(t)}\|_F^2=\|{U(t)^H\dot U_j(t)}\|_F^2=\sum_{i=1,i\ne j}^k\|{U_i(t)^H\dot U_j(t)}\|_F^2.
\end{eqnarray}
Then we get
\begin{eqnarray}\label{c-2-10}
\|{\Delta A}\|_F^2&=&\|{U(t)^H\Delta AU(t)}\|_F^2=\sum_{j=1}^k\|{\Delta A_{jj}(t)}\|_F^2
+\sum_{i\ne j}\|{\Delta A_{ij}(t)}\|_F^2\nonumber\\
&\ge &\sum_{j=1}^k\|{\dot\Lambda_j(t)}\|_F^2
+\sum_{j=1}^k\sum_{i=1,i\ne j}^k\delta_{ji}(t)^2\|{U_i(t)^H\dot U_j(t)}\|_F^2   \qquad (\mbox{by}\; (\ref{c-2-5})\; \mbox{and}\; (\ref{c-2-8}) )\nonumber\\
&\ge &\|{\dot\Lambda(t)}\|_F^2+\sum_{j=1}^k\delta_j(t)^2\sum_{i=1,i\ne j}^k\|{U_i(t)^H\dot U_j(t)}\|_F^2  \qquad (\mbox{by}\; (\ref{c-2-1}))\nonumber\\
&=&\|{\dot\Lambda(t)}\|_F^2
+\sum_{j=1}^k\delta_j(t)^2\|{\dot U_j(t)}\|_F^2\qquad (\mbox{by}\; (\ref{c-2-9}))
\nonumber\\
&\ge& \|{\dot\Lambda(t)}\|_F^2+\sum_{j=1}^k\delta_{j,\min}^2\|{\dot U_j(t)}\|_F^2.\qquad (\mbox{by}\; (\ref{c-2-2}))
\end{eqnarray}
Using the fact that
\[
\left|\int_0^1f(t){\rm d}t\right|^2\leq\int_0^1|f(t)|^2{\rm d}t,\quad \forall f(t),
\]
and together with (\ref{c-2-10}), we get
\begin{eqnarray}\label{c-2-11}
\|{\Delta A}\|_F^2&=&\int_0^1\|{\Delta A}\|_F^2{\rm d}t\geq
\int_0^1\|{\dot\Lambda(t)}\|_F^2{\rm d}t
+\sum_{j=1}^k\delta_{j,\min}^2\int_0^1\|{\dot U_j(t)}\|_F^2{\rm d}t\nonumber\\
&\geq&
\left\|\int_0^1\dot\Lambda(t){\rm d}t\right\|_F^2
+\sum_{j=1}^k\delta_{j,\min}^2\left\|\int_0^1\dot U_j(t){\rm d}t\right\|_F^2\nonumber\\
&=&\|{\Lambda(1)-\Lambda(0)}\|_F^2+\sum_{j=1}^k\delta_{j,\min}^2\|{U_j(1)-U_j(0)}\|_F^2
\nonumber\\
&=&\|{\Lambda(1)-\Lambda(0)}\|_F^2+\sum_{j=1}^k\delta_{j,\min}^2\|{\wt U_j-U_j}\|_F^2.
\end{eqnarray}
By (\ref{c-2-11}) and the inequality
\[
\|{\rm Eig}^{\downarrow}(\wt A)-{\rm Eig}^{\downarrow}(A)\|_F=\|{\rm Eig}^{\downarrow}(\Lambda(1))-{\rm Eig}^{\downarrow}(\Lambda(0))\|_F\leq\|{\Lambda(1)-\Lambda(0)}\|_F,
\]
which is from (\ref{c-121}),  we obtain (\ref{c-2-4}).
\end{proof}

\end{document}
